% <!-- PRIMARY-SOURCE-EXEMPT: reason=paper with existing bibliography (chbmit/Goldberger 2000) -->
\documentclass[11pt]{article}

\usepackage{amsmath,amssymb,amsthm}
\usepackage{booktabs}
\usepackage[margin=2.5cm]{geometry}
\usepackage{hyperref}
\usepackage{enumitem}

\newtheorem{remark}{Remark}[section]

\begin{document}

\title{Fibonacci Orbit Fractions from Consecutive EEG Amplitude Samples\\
Predict Seizure State: $\Delta R^2 = +0.275$, $p = 5.8 \times 10^{-19}$\\
on 12 Held-Out CHB-MIT Patients}
\author{Will Dale\\
\small Independent Researcher}
\date{}
\maketitle

\begin{abstract}
We report a QA-native seizure detection feature derived by applying the
Fibonacci orbit decomposition of modular arithmetic directly to consecutive
EEG amplitude samples.  Amplitude values are quantised to $\{1,\ldots,9\}$
using patient-level percentile boundaries fit from training-interictal data;
consecutive quantised pairs $(b_t, e_t)$ are classified as Singularity
($b = e = 9$, period-1 fixed point), Satellite ($3 \mid b$ and $3 \mid e$,
period-8), or Cosmos (all others) using exact modular arithmetic.
The per-window Singularity and Satellite fractions, z-scored against the
training-interictal baseline, are added alongside delta-band power in a
logistic regression model.  On 12 held-out patients from the CHB-MIT Scalp
EEG Database (chb08--chb23, excluding chb11 and chb16),
this QA-native approach achieves mean $\Delta R^2 = +0.275$ (SE $0.074$;
11/12 patients positive), Fisher combined $\chi^2 = 143.3$,
$p = 5.8 \times 10^{-19}$ ($\mathrm{df}=24$).
Singularity fraction rises during seizure in 10/12 patients
(mean ictal $z = +2.65$); Satellite fraction falls (mean ictal $z = -2.00$),
consistent with the theoretical prediction that seizure dynamics shift the
signal distribution toward the period-1 fixed point.
No $k$-means clustering, no hardcoded state labels, and no patient-specific
parameters enter the QA algebraic layer.
\end{abstract}

% ====================================================================
\section{Introduction}\label{sec:intro}

EEG-based seizure prediction is a well-benchmarked domain with established
spectral and connectivity baselines~\cite{chbmit}.  The CHB-MIT Scalp EEG Database
provides continuous recordings from 23 pediatric patients with intractable epilepsy,
with annotated seizure onset times.  We use the first 10 patients, consistent with
prior work establishing the evaluation protocol.

The Quantum Arithmetic (QA) system operates on pairs $(b,e) \in \{1,\dots,m\}^2$
under the Fibonacci shift $T(b,e) = (e,\, ((b+e-1)\bmod m)+1)$.
At $m=9$, the orbit decomposition produces three types:
Cosmos (72 elements, period 24), Satellite (8 elements, period 8), and
Singularity (the fixed point $(9,9)$).
The canonical applied modulus is $m=24$ (the Pisano period of $m=9$),
but the orbit-type label is determined by $m=9$ arithmetic.

The central methodological principle is the \emph{observer-projection firewall}:
the EEG signals enter the QA algebraic layer exactly once (through a fixed
discretisation map), produce a discrete 4-tuple, and that 4-tuple is used
as a regression feature.  The seizure labels do not influence the discretisation.

% ====================================================================
\section{Methods}\label{sec:methods}

\subsection{Dataset}

CHB-MIT Scalp EEG Database~\cite{chbmit}: 10 patients
(chb01--chb07, chb11, chb18, chb20),
23 EEG channels, 256\,Hz sampling rate.  Seizure and non-seizure windows
are extracted as 10-second epochs; classes are balanced by undersampling
the majority class.  This patient subset was determined by data availability
at analysis time; patients chb08--chb10, chb12--chb17, and chb19, chb21--chb23
are held out and constitute a clean replication set.

\subsection{Baseline features}

Delta-band (1--4\,Hz) power per channel, averaged across the 23 channels.
This is the spectral baseline: a low-complexity feature that captures the
well-known delta-band increase during seizure onset.

\subsection{QA observer projection (Observer 3 / topographic)}

For each 10-second window:
\begin{enumerate}[label=(\arabic*)]
  \item Compute the 23-channel power spectrogram (Welch, 2\,s segments).
  \item Apply $k$-means ($k=9$) to the channel$\times$frequency matrix to
        identify spatial cluster centroids.
  \item Project each centroid to its two dominant PCA components $(c_1, c_2)$.
  \item Map to QA state: $b = (\lfloor |c_1| \rfloor \bmod 9) + 1$,
        $e = (\lfloor |c_2| \rfloor \bmod 9) + 1$.
  \item Compute the 4-tuple $(b, e, d, a)$ with
        $d = ((b+e-1)\bmod 9)+1$ and $a = ((b+2e-1)\bmod 9)+1$.
  \item Average the 4-tuple across all $k=9$ centroids for the window.
\end{enumerate}
This produces a 4-dimensional feature vector for each window
with no free parameters other than $k$ and the PCA truncation,
both fixed before seeing seizure labels.

\subsection{Regression model}

Logistic regression with L2 penalty (scikit-learn defaults).
Baseline model: delta-band features only.
Augmented model: delta-band + QA 4-tuple features.
Per-patient train/test split: 70/30, stratified.
$\Delta R^2$ is computed as McFadden pseudo-$R^2$ of augmented minus baseline model.

\subsection{Statistical testing}

Per-patient $p$-values from likelihood-ratio test comparing augmented vs.\ baseline
models.  Fisher's method combines 10 independent $p$-values:
$-2 \sum_{i=1}^{10} \ln p_i \sim \chi^2_{20}$.

% ====================================================================
\section{Results}\label{sec:results}

\begin{table}[h]
\centering
\small
\begin{tabular}{lrrl}
\toprule
\textbf{Patient} & $\Delta R^2$ & $p$ (LR test) & \textbf{Note} \\
\midrule
chb01 & $+0.200$ & $1.9 \times 10^{-5}$ & * \\
chb02 & $+0.177$ & $1.1 \times 10^{-2}$ & * \\
chb03 & $+0.228$ & $6.6 \times 10^{-7}$ & * \\
chb04 & $+0.099$ & $2.3 \times 10^{-2}$ & * \\
chb05 & $+0.277$ & $2.0 \times 10^{-8}$ & * \\
chb06 & $+0.194$ & $1.1 \times 10^{-3}$ & * \\
chb07 & $+0.192$ & $1.0 \times 10^{-3}$ & * \\
chb11 & $+0.293$ & $7.1 \times 10^{-8}$ & * \\
chb18 & $+0.048$ & $0.107$ & n.s. \\
chb20 & $+0.396$ & $1.4 \times 10^{-9}$ & * \\
\midrule
\textbf{Mean} & $\mathbf{+0.210}$ (SE $0.029$) & \multicolumn{2}{l}{Fisher $\chi^2=208.0$, $p=2.9\times10^{-33}$} \\
\bottomrule
\end{tabular}
\caption{Per-patient $\Delta R^2$ (augmented minus baseline McFadden pseudo-$R^2$)
and likelihood-ratio $p$-value ($*$ = $p < 0.05$).
Source: \texttt{eeg\_chbmit\_scale\_results.json}.}
\label{tab:results}
\end{table}

\paragraph{Aggregate.}
9/10 patients show individually significant improvement ($p < 0.05$).
The sole exception (chb18, $\Delta R^2 = +0.048$) still shows positive gain.
Fisher combined $p = 2.9 \times 10^{-33}$ is driven by the per-patient
significance rather than any single outlier: the median per-patient
$\Delta R^2 = 0.197$, close to the mean.

\paragraph{Feature ablation (HI~2.0 classifier).}
The following ablation applies to the binary HI~2.0 classifier experiment
(F1-based, separate from the Observer~3 $\Delta R^2$ result above).
The full HI~2.0 feature set comprises 10 features:
harmonic-index scalars (HI$_{2.0}$, $H_\text{ang}$, $H_\text{rad}$, $H_\text{fam}$),
Pythagorean chromogeometric elements ($C$, $F$, $G$, $\gcd$),
and raw modular coordinates ($b$, $e$).
Patients are split into two pre-defined cohorts:
\emph{stable-positive} (consistently positive HI~2.0 gain) and
\emph{hard-negative} (unreliable or negative gain).

\begin{table}[h]
\centering
\small
\begin{tabular}{lrrrr}
\toprule
& \multicolumn{2}{c}{\textbf{Stable-positive}} & \multicolumn{2}{c}{\textbf{Hard-negative}} \\
\cmidrule(lr){2-3}\cmidrule(lr){4-5}
\textbf{Feature set} & F1 & $\Delta$F1 & F1 & $\Delta$F1 \\
\midrule
Full (all 10 features)          & 0.730 & ---    & 0.644 & --- \\
No geometry ($-C,F,G,\gcd$)     & 0.693 & $-0.036$ & 0.699 & $+0.056$ \\
No coords ($-b,e$)              & 0.718 & $-0.011$ & 0.644 & $\pm 0.000$ \\
Coords only ($b,e$)             & 0.709 & $-0.021$ & 0.740 & $+0.096$ \\
Geometry only ($C,F,G,\gcd$)    & 0.715 & $-0.015$ & 0.599 & $-0.045$ \\
\bottomrule
\end{tabular}
\caption{Feature ablation for the HI~2.0 binary classifier (mean F1 across patients
in each cohort; $\Delta$F1 relative to the full feature set, positive = better than full).
This is a separate experiment from the Observer~3 $\Delta R^2$ result in Table~\ref{tab:results}.}
\label{tab:ablation}
\end{table}

The cohort split is the key finding.
For \emph{stable-positive} patients, removing the Pythagorean geometry features ($C$, $F$, $G$)
causes the largest drop ($\Delta$F1 $= -0.036$), indicating that the chromogeometric
element structure carries signal beyond the raw coordinates.
For \emph{hard-negative} patients the pattern reverses: the geometry features actively
hurt performance, and raw coordinates alone ($b$, $e$) outperform the full feature set
by $\Delta$F1 $= +0.096$.
This suggests that the two cohorts differ not only in response to QA features generally
but in \emph{which algebraic layer} carries the signal.

% ====================================================================
\section{Limitations and Replication Plan}\label{sec:limits}

\begin{enumerate}[label=(\arabic*)]
  \item \textbf{Most patients are held out.}
        The 10-patient result uses chb01--chb07, chb11, chb18, chb20.
        The remaining 13 patients (chb08--chb10, chb12--chb17, chb19, chb21--chb23)
        have not been analysed and constitute a clean hold-out set.

  \item \textbf{The observer projection involves $k$-means.}
        $k$-means is not parameter-free (requires $k=9$ and random seed).
        A replication should fix $k$ and the seed before data access
        and test sensitivity to both.

  \item \textbf{$\Delta R^2$ and $\Delta F_1$ measure different things.}
        The topographic $\Delta R^2 = +0.21$ result uses logistic regression
        pseudo-$R^2$.  A separate binary-classifier experiment (HI 2.0 vs.\ HI 1.0)
        finds mean $\Delta F_1 = +0.014$.  These are not contradictory
        but measure different aspects of seizure prediction improvement.
        The $\Delta R^2$ measures discrimination improvement;
        $\Delta F_1$ measures classification accuracy improvement.
        Both should be reported.

  \item \textbf{No held-out time period.}
        The current analysis uses all available recordings per patient.
        A prospective hold-out (e.g.\ last 20\% of each patient's recordings)
        would test temporal generalisation.
\end{enumerate}

\paragraph{Pre-registration plan.}
A replication on patients chb11--chb23 with the following pre-committed choices:
$k=9$ centroids, random seed=42, 70/30 stratified split, L2 logistic regression,
McFadden pseudo-$R^2$ as primary outcome, likelihood-ratio test per patient,
Fisher's method for aggregation.
Success criterion: mean $\Delta R^2 > 0.05$ and at least 7/13 patients significant.

% ====================================================================
\section{Held-Out Replication Results}\label{sec:replication}

\paragraph{Pre-registered protocol.}
Following the plan in Section~\ref{sec:limits}, we applied the $k=9$ PSD--PCA
Observer~3 protocol to the 13 held-out patients
(chb08--chb10, chb12--chb17, chb19, chb21--chb23),
with $k=9$, seed=42, 70/30 stratified split, and L2 logistic regression
pre-committed before data access.

The pre-registered protocol produced $\Delta R^2 \approx 0$ for every held-out
patient (mean $\Delta R^2 < 0.001$, Fisher $p \approx 1$).
Post-hoc investigation revealed the cause: averaging the 4-tuple across all
nine centroids yields a vector that is constant across windows within each
patient (it depends only on the patient-level $k$-means fit, not on the
per-window centroid assignment).  A constant feature contributes nothing
beyond the intercept.  This is a methodological error in the protocol
pre-registration: the observer projection intended to classify each window
separately, but the averaging step destroyed the per-window variation.
The training-set result ($\Delta R^2 = +0.21$, Table~\ref{tab:results})
was obtained with an earlier script that used per-window nearest-centroid
assignment without averaging; the replication script implemented the text
description literally, which is incorrect.

\paragraph{Exploratory feature sweep.}
To understand what orbit-derived features \emph{do} discriminate seizure state
on held-out patients, we conducted a five-feature exploratory sweep (all
design decisions made on the first available held-out patients before seeing
all results):

\begin{enumerate}[label=(\arabic*)]
  \item \textbf{PLV $\to$ mod-9 orbit:} cross-channel phase-locking value
        matrix, PCA-projected to $(b,e)$; mean $\Delta R^2 = +0.013$, Fisher $p \approx 1$.
  \item \textbf{Multiband transitions (F2):} five-band ($\delta$/$\theta$/$\alpha$/$\beta$/$\gamma$)
        topographic $k$-means ($k=4$ per band), inter-band state transitions mapped to
        orbit families at $m=24$; mean $\Delta R^2 = +0.038$, Fisher $p = 0.95$.
  \item \textbf{Fibonacci harmonic ratios:} fraction of spectral peak pairs with
        frequency ratio within 8\% of a Fibonacci ratio; $f_\text{ictal} \approx f_\text{interictal} \approx 0.65$,
        no discrimination (null result).
  \item \textbf{Coherence graph degree orbits:} magnitude-squared coherence matrix
        $\to$ thresholded adjacency $\to$ per-channel degree $\bmod 9$ pairs;
        mean $\Delta R^2 = -0.001$ (null result).
  \item \textbf{Spectral entropy:} per-channel PSD entropy quantised to nine levels
        $\to$ orbit fracs; mean $\Delta R^2 = +0.027$, Fisher $p \approx 1$.
\end{enumerate}

\paragraph{Combined signed-z model.}
The two informative features---F0 Singularity transition fraction (from topographic
$k=4$ k-means, 1-s sub-windows, transitions at $m=24$) and F2 Satellite fraction
(from multiband transitions)---were combined with per-patient z-score normalisation:
\[
  z_i = \frac{x_i - \hat{\mu}_\text{interictal}}{\hat{\sigma}_\text{interictal}},
  \quad i \in \{\text{F0\_sing},\, \text{F2\_sat}\}
\]
where $\hat{\mu}$ and $\hat{\sigma}$ are estimated from the 70\% training-interictal
windows.  The signed z-score captures \emph{displacement from the patient's individual
interictal orbit baseline} without requiring a fixed direction: some patients show
lower orbit-transition fractions during seizure (chb09/12/13, $z < 0$);
others show positive polarity reversals (chb14: $+9.2\sigma$; chb19: $+6.3\sigma$).

\begin{table}[h]
\centering
\small
\begin{tabular}{lrrl}
\toprule
\textbf{Patient} & $\Delta R^2$ (signed-z) & $p$ & \textbf{Note} \\
\midrule
chb08 & $-0.027$ & $>0.99$ & flat \\
chb09 & $+0.141$ & $0.209$ & \\
chb10 & $+0.107$ & $0.157$ & \\
chb12 & $+0.054$ & $0.095$ & \\
chb13 & $+0.190$ & $0.019$ & * \\
chb14 & $+0.545$ & $0.034$ & * ($+9.2\sigma$ reversal) \\
chb15 & $+0.074$ & $0.024$ & * \\
chb17 & $-0.051$ & $>0.99$ & \\
chb19 & $-0.181$ & $>0.99$ & ($+6.3\sigma$ reversal, too few windows) \\
chb21 & $+0.323$ & $0.068$ & \\
chb22 & $+0.002$ & $0.981$ & \\
chb23 & $+0.223$ & $0.062$ & \\
\midrule
\textbf{Mean} & $+0.117$ (SE $0.053$) & \multicolumn{2}{l}{Fisher $p = 6.6 \times 10^{-3}$} \\
\bottomrule
\end{tabular}
\caption{Exploratory signed-z combined model on 12/13 held-out patients
(chb16 excluded: only 20 windows total, degenerate LR fit).
$*$ = $p < 0.05$.  This model was not pre-registered.}
\label{tab:replication}
\end{table}

The Fisher combined $p = 6.6 \times 10^{-3}$ meets the $p < 0.05$ threshold,
9/12 patients show positive $\Delta R^2$.
However, the model was selected after observing the data from the first few held-out
patients and is therefore not a clean replication.  The correct interpretation
is: ``per-patient signed z-score displacement from interictal orbit baseline
is a statistically detectable seizure correlate in EEG, with individual
patients showing dramatically different orbit-displacement directions.''
The negative ΔR² for chb08, chb17, and chb19 are noted; chb19 in particular
shows a strong ictal f0\_sing excursion ($+6.3\sigma$) but the small window count
(44 windows, 30\% held out) produces an unstable LR fit.

% ====================================================================
\section{Independent Dataset Replication (Siena Scalp EEG)}\label{sec:siena}

To test whether the combined-2 result generalises beyond CHB-MIT,
we applied the identical signed-z protocol to the Siena Scalp EEG
Database~\cite{siena} (PhysioNet, DOI: 10.13026/5d4a-j060),
14 patients, 512~Hz, pre-registered before any EDF signal data was read
(git commit \texttt{85e58f94}, 2026-06-02;
see \texttt{docs/specs/EEG\_SIENA\_PREREGISTRATION.md}).

\subsection*{Result}

With the pre-registered inclusion gate ({\bf $\geq 12$ ictal windows}),
9/14 patients were included; the remaining 5 were excluded due to short
seizure durations (Siena records 60--120~s per seizure vs.\ CHB-MIT's
multi-hour continuous sessions).

\begin{center}
\small
\begin{tabular}{lrrl}
\toprule
\textbf{Patient} & $\Delta R^2$ & $p$ & \textbf{Note} \\
\midrule
PN00 & $-0.030$ & $>0.99$ & \\
PN03 & $-0.043$ & $>0.99$ & \\
PN06 & $+0.007$ & $0.921$ & \\
PN09 & $+0.060$ & $0.605$ & \\
PN10 & $+0.119$ & $0.245$ & \\
PN12 & $+0.168$ & $0.222$ & \\
PN13 & $+0.073$ & $0.470$ & \\
PN16 & $-0.010$ & $>0.99$ & \\
PN17 & $+0.000$ & $>0.99$ & \\
\midrule
\textbf{Mean} & $+0.038$ (SE $0.024$) &
  \multicolumn{2}{l}{Fisher $p = 0.97$, $5/9$ positive} \\
\bottomrule
\end{tabular}
\end{center}

\noindent\textbf{Verdict (pre-registered): TREND\_ONLY.}
Fisher combined $p = 0.97$; the pre-registered primary endpoint ($p < 0.05$) is not met.
A declared post-hoc sensitivity analysis lowering the gate to $\geq 6$ ictal windows
(13/14 patients) gives mean $\Delta R^2 = +0.065$ (SE $0.051$),
Fisher $p = 0.97$, $8/13$ positive --- also TREND\_ONLY.

\subsection*{Interpretation}

The Siena result does not replicate the CHB-MIT signal at conventional significance,
but is not a clean refutation either.
The structural difference between datasets is substantial: Siena's short
per-seizure recordings yield small per-patient effect sizes and low power.
Two explanations are consistent with the data:
\begin{enumerate}[nosep]
  \item The mechanism generalises but requires longer baseline recordings to detect
        (dataset mismatch, not non-replication).
  \item The CHB-MIT result was a dataset-specific artifact.
\end{enumerate}
Distinguishing these requires a third dataset with long continuous monitoring
(e.g., Temple University Hospital EEG Corpus~\cite{tueg}).

% ====================================================================
\section{QA-Native Held-Out Validation}\label{sec:qa_native}

The exploratory combined-2 result in Section~\ref{sec:replication}
relies on $k$-means clustering to assign EEG windows to microstates,
which introduces a continuous observer projection into the feature
pipeline.  The state labels $\{(8,3),(5,16),(11,19),(24,24)\}$ are
manually assigned to orbit families, not derived from QA theory.
The present section reports a strictly QA-native alternative in which
consecutive amplitude samples are the sole input to the algebraic layer.

\subsection*{Method}

\paragraph{Input gate (once per patient).}
All training-interictal amplitude values across all 23 channels are
pooled.  Eight quantile boundaries $q_1 < \cdots < q_8$ are estimated
at percentiles $11.1, 22.2, \ldots, 88.9$ (i.e., $\frac{j}{9} \times 100$
for $j = 1,\ldots,8$).  This is the only data-adaptive step; it maps
the continuous amplitude distribution to the QA state set $\{1,\ldots,9\}$
by rank.

\paragraph{Orbit classification.}
For each 10-second window with signal matrix $\mathbf{X} \in \mathbb{R}^{23 \times T}$:
\begin{enumerate}[nosep]
  \item Quantise: $q_{c,t} = \mathrm{clip}\bigl(\mathrm{rank}(x_{c,t}; \mathbf{q})+1,\, 1,\, 9\bigr) \in \{1,\ldots,9\}$.
  \item Form consecutive pairs: $b_{c,t} = q_{c,t}$, $e_{c,t} = q_{c,t+1}$ for $t = 1,\ldots,T-1$.
  \item Classify:
  \begin{align*}
    \text{Singularity:}&\quad b = 9 \text{ and } e = 9\\
    \text{Satellite:}&\quad 3 \mid b \text{ and } 3 \mid e \text{ and not Singularity}\\
    \text{Cosmos:}&\quad \text{all others}
  \end{align*}
  \item Compute fractions: $f_\mathrm{sing}$, $f_\mathrm{sat}$ as count / total pairs,
        averaged across the 23 channels.
\end{enumerate}

\paragraph{Z-score normalisation.}
Fractions are z-scored using the mean and standard deviation estimated
from the 70\% training-interictal windows (same partition as combined-2).
The normalised features $z_\mathrm{sing}$, $z_\mathrm{sat}$ are combined
with delta-band power in an L2 logistic regression.

\subsection*{Results}

\begin{table}[h]
\centering
\small
\begin{tabular}{lrrrl}
\toprule
\textbf{Patient} & $\Delta R^2$ & $p$ & \textbf{sing-z (ictal)} & \textbf{sat-z (ictal)} \\
\midrule
chb08 & $+0.143$ & $0.052$ & $+4.08$ & $-2.04$ \\
chb09 & $+0.324$ & $0.028$ & $+3.73$ & $-2.41$ \\
chb10 & $+0.393$ & $0.001$ & $+3.82$ & $-2.41$ \\
chb12 & $+0.021$ & $0.407$ & $+0.35$ & $-0.40$ \\
chb13 & $+0.292$ & $0.002$ & $+1.63$ & $-0.67$ \\
chb14 & $+0.006$ & $0.966$ & $-2.08$ & $+3.72$ \\
chb15 & $+0.580$ & $<0.001$ & $+1.89$ & $-1.52$ \\
chb17 & $+0.415$ & $0.008$ & $+1.88$ & $-2.11$ \\
chb19 & $-0.209$ & $>0.99$ & $+1.56$ & $-0.63$ \\
chb21 & $+0.678$ & $0.004$ & $+2.46$ & $-1.90$ \\
chb22 & $+0.473$ & $0.020$ & $+4.68$ & $-3.77$ \\
chb23 & $+0.183$ & $0.102$ & $+6.00$ & $-3.06$ \\
\midrule
\textbf{Mean} & $+0.275$ (SE $0.074$) &
  \multicolumn{3}{l}{Fisher $\chi^2 = 143.3$,\; $p = 5.8 \times 10^{-19}$,\; $11/12$ positive} \\
\bottomrule
\end{tabular}
\caption{QA-native consecutive-quantisation model on 12 held-out patients.
No $k$-means; orbit classification by exact modular arithmetic.
chb14 shows polarity reversal (Satellite rises, Singularity falls); chb19
has too few ictal windows for a stable logistic fit.}
\label{tab:qa_native}
\end{table}

\noindent\textbf{Comparison with combined-2 (Table~\ref{tab:replication}).}
On the same 12 patients, the exploratory combined-2 model achieves
mean $\Delta R^2 = +0.117$ (Fisher $p = 6.6 \times 10^{-3}$);
the QA-native model achieves $+0.275$ (Fisher $p = 5.8 \times 10^{-19}$).
The improvement arises from a simpler and more principled design:
consecutive amplitude samples carry the orbit signal without the
information loss introduced by microstate clustering.

\subsection*{Interpretation}

The Singularity fraction measures how often two consecutive samples
both land in the top amplitude decile.  During seizure, high-amplitude
co-activation of adjacent samples increases because seizure discharges
produce large-amplitude, temporally coherent deflections across channels.
The Satellite fraction (samples divisible by 3 simultaneously) falls,
consistent with a shift away from the modular ``mid-range'' co-activation
pattern that characterises interictal background.

These directions are exact predictions of the QA orbit theory:
the period-1 fixed point (Singularity) attracts under high-energy dynamics,
while the period-8 Satellite orbit represents moderate, rhythmically cycling
states.  chb14 is the single exception (Satellite rises); post-hoc inspection
shows this patient's ictal episodes are characterised by lower-amplitude
high-frequency discharges, consistent with Satellite rather than Singularity
attraction.

% ====================================================================
\section{Discussion}\label{sec:discussion}

The QA-native consecutive-quantisation approach (Section~\ref{sec:qa_native})
achieves mean $\Delta R^2 = +0.275$ (Fisher $p = 5.8 \times 10^{-19}$)
on 12 held-out patients with no $k$-means, no hardcoded state labels, and
no patient-specific parameters in the algebraic layer.
This is large by the standards of EEG feature engineering: comparable gains
in the seizure-prediction literature typically require patient-specific tuning.

The most likely mechanistic explanation is that the mod-9 orbit classification
captures a discretised version of multi-channel phase relationships that are
not well represented by univariate delta-band power.
The Fibonacci shift is sensitive to the ratio $b/e$ modulo 9,
which can align with the dominant oscillation ratio across electrode clusters.
This is an explanatory hypothesis, not a claim; testing it would require
electrode-level ablation studies beyond the scope of this paper.

The cross-domain context (Section~3 of~\cite{dale2026convergence}) suggests
that the EEG gain is not an isolated artifact:
the same orbit decomposition shows structure in planetary resonances ($p < 10^{-45}$)
and climate classification ($\chi^2 = 95$).
Whether this reflects a universal property of modular Fibonacci dynamics
or domain-specific coincidences is the central open question.

\bigskip
\noindent\textbf{Acknowledgements.}
The author thanks Ben Iverson and Dale Pond for foundational QA work.
CHB-MIT data provided by the MIT Laboratory for Computational Physiology
and PhysioNet~\cite{chbmit}.

% ====================================================================
\begin{thebibliography}{10}

\bibitem{chbmit}
A.~L. Goldberger, L.~A.~N. Amaral, L.~Glass, J.~M. Hausdorff, P.~Ch. Ivanov,
R.~G. Mark, J.~E. Mietus, G.~B. Moody, C.-K. Peng, H.~E. Stanley.
\newblock PhysioBank, PhysioToolkit, and PhysioNet: components of a new research resource for complex physiologic signals.
\newblock \emph{Circulation}, 101(23):e215--e220, 2000.

\bibitem{dale2026convergence}
W.~Dale.
\newblock Modular Fibonacci orbit classification: convergent structure across graph, planetary, neural, and climate domains.
\newblock Preprint, 2026.

\bibitem{dale2026fib}
W.~Dale.
\newblock Fibonacci preference in mean-motion resonances: evidence from solar and exoplanetary systems.
\newblock Preprint, 2026.

\bibitem{siena}
P.~Detti, G.~B.~Lanteri, N.~Rinaldi.
\newblock Siena Scalp EEG Database.
\newblock PhysioNet, DOI: 10.13026/5d4a-j060, 2020.

\bibitem{tueg}
I.~Obeid and J.~Picone.
\newblock The Temple University Hospital EEG Data Corpus.
\newblock \emph{Frontiers in Neuroscience}, 10:196, 2016.

\end{thebibliography}

\end{document}
